\section{Metoda}
\label{sec:metoda}

W tej sekcji opisujemy szczegóły algorytmiczne i implementacyjne wszystkich
agentów porównywanych w eksperymentach.

\subsection{Reprezentacja stanu gry}

Stan gry Breakthrough w naszej implementacji jest reprezentowany przez parę
bitboardów $u_{64}$: jeden dla pionków białych, jeden dla czarnych. Bit $i$
odpowiada polu $i = r \cdot c + k$, gdzie $r$ to numer rzędu (0-indeksowany
od dołu), $k$ to numer kolumny, a $c$ to liczba kolumn planszy. Dla planszy
6$\times$6 wszystkie 36 pól mieszczą się w jednym \texttt{u64}, dla 8$\times$8
zajmujemy całe 64 bity. Stan gry zawiera również:
flagę \texttt{white\_to\_move};
licznik \texttt{move\_count};
wartość hash Zobrist $\in \{0, \ldots, 2^{64}-1\}$ zaktualizowaną przyrostowo;
stałe maski (lewa i prawa krawędź, plansza, rzędy zwycięskie).

\paragraph{Generowanie ruchów.}
Korzystamy z operacji bitowych z maskowaniem krawędzi przed przesunięciem,
co eliminuje błąd ,,wraparound''. Dla białego (porusza się ku wyższym wierszom):
\begin{align*}
\text{prosto} &= (W \ll c) \,\&\, \text{empty} \,\&\, \text{board\_mask} \\
\text{ukos lewy} &= ((W \,\&\, \overline{\text{left\_edge}}) \ll (c-1)) \,\&\, (\text{empty} \cup B) \,\&\, \text{board\_mask} \\
\text{ukos prawy} &= ((W \,\&\, \overline{\text{right\_edge}}) \ll (c+1)) \,\&\, (\text{empty} \cup B) \,\&\, \text{board\_mask}
\end{align*}
gdzie $W$, $B$ to bitboardy białych i czarnych, $\overline{X}$ to negacja.
Dla czarnego (przesunięcia w dół) ruchy są symetryczne, używają operatora $\gg$.

\paragraph{Detekcja końca gry.}
Cztery warunki sprawdzane w kolejności (od najtańszego do najdroższego):
(1) biały dotarł do ostatniego rzędu;
(2) czarny dotarł do rzędu 0;
(3) jeden z graczy stracił wszystkie pionki;
(4) gracz na ruchu nie ma legalnych ruchów (przegrywa).

\paragraph{Hash Zobrist.}
Tablica losowych $u_{64}$ rozmiar $2 \cdot \text{rows} \cdot \text{cols}$
generowana raz, deterministycznie seedowana. Hash stanu = XOR wartości tablicy
dla każdego pionka na planszy. Współdzielona przez referencję
(\texttt{Arc<Vec<u64>>}) między wszystkimi instancjami stanu, dzięki czemu klonowanie
\texttt{GameState} jest stałoczasowe.

\paragraph{Walidacja krzyżowa.}
Dla pewności poprawności silnika, równolegle z implementacją w Ruście została
napisana referencyjna implementacja w czystym Pythonie (priorytet czytelności).
Test \texttt{test\_rust\_vs\_python.py} gra 1000 losowych partii i porównuje
po każdym ruchu zbiory legalnych ruchów, stan terminalny, zwycięzcę i
reprezentację tekstową. Pełna zgodność jest warunkiem koniecznym przed
uruchomieniem eksperymentów.

\subsection{Vanilla UCT (Kocsis \& Szepesvári 2006)}

Klasyczny algorytm 4-fazowego MCTS \cite{kocsis2006bandit}:

\begin{enumerate}[nosep]
    \item \textbf{Selekcja.} Począwszy od korzenia, w każdym wewnętrznym węźle
          $s$ wybieramy potomka maksymalizującego UCB1:
          \begin{equation}
              \text{score}(c) = \frac{\text{wins}(c)}{\text{visits}(c)}
                              + c_{\text{exp}} \sqrt{\frac{\ln \text{visits}(s)}{\text{visits}(c)}}
              \label{eq:ucb1}
          \end{equation}
          Dla potomka z $\text{visits} = 0$ zwracamy $\infty$ (forsujemy
          jego eksplorację).
    \item \textbf{Ekspansja.} W węźle z nierozwiniętymi ruchami wybieramy
          jeden losowy nierozwinięty ruch i tworzymy nowy węzeł potomny.
    \item \textbf{Symulacja.} Od nowego węzła gramy losowo (uniform) do
          terminala.
    \item \textbf{Backpropagation.} Wynik symulacji przekazywany jest
          w górę drzewa; na każdym węźle aktualizujemy:
          $\text{visits}(n) \mathrel{{+}{=}} 1$, oraz
          $\text{wins}(n) \mathrel{{+}{=}} 1$ jeśli wygrał gracz, który zagrał
          ruch \emph{do} tego węzła.
\end{enumerate}

\paragraph{Konwencja perspektywy.}
$\text{wins}(c)$ liczone jest z perspektywy gracza, który wykonał ruch $c$
(czyli gracza znajdującego się na ruchu w rodzicu $c$). Ta konwencja
zapewnia, że UCB1 w wzorze (\ref{eq:ucb1}) wybiera ruch dający
maksymalne wygrane dla gracza wybierającego --- dokładnie to, czego oczekujemy
od minimax-spójnej selekcji.

\paragraph{Implementacja.}
Pula węzłów to płaska \texttt{Vec<UctNode>} indeksowana \texttt{usize}
zamiast wskaźników. Węzły nie przechowują pełnego \texttt{GameState} (oszczędność
pamięci); stan przy schodzeniu od korzenia jest rekonstruowany przez
ponowne aplikowanie ruchów. Cena: $O(\text{depth})$ klonowań stanu na
iterację (akceptowalne ze względu na taniość klonowania $\sim$80B + Arc).

Eksponowany do Pythona przez PyO3 jako \texttt{UctAgent(iterations, c, seed)}.

\subsection{UCT + RAVE (Gelly \& Silver 2011)}

\paragraph{Algorytm.}
Każdy węzeł $s$ przechowuje, oprócz statystyk UCT, dodatkowe statystyki AMAF
(All-Moves-As-First): mapę kluczowaną kodem ruchu \texttt{from*64 + to},
zawierającą liczbę odwiedzin $\tilde{n}(s, a)$ i wygranych $\tilde{w}(s, a)$
dla ruchu $a$ jako pierwszego z $s$.

W fazie selekcji potomek $c$ rodzica $s$ jest oceniany formułą zmieszaną
\cite{gelly2011monte}:
\begin{align}
    \text{score}(c) &= (1 - \beta) Q_{\text{UCT}}(c) + \beta Q_{\text{RAVE}}(c) + c_{\text{exp}} \sqrt{\frac{\ln N(s)}{n(c)}} \label{eq:rave}\\
    \beta(n, \tilde{n}) &= \frac{\tilde{n}}{n + \tilde{n} + 4 n \tilde{n} b^2} \label{eq:beta}
\end{align}
gdzie $Q_{\text{UCT}}(c) = \text{wins}(c) / \text{visits}(c)$ to ocena UCT,
$Q_{\text{RAVE}}(c) = \tilde{w}(s, a_c) / \tilde{n}(s, a_c)$ to wartość AMAF
ruchu $a_c$ wiodącego do $c$ (zapisana w $s$, indeksowana kodem ruchu),
$b^2$ to hiperparametr formuły Silvera (\ref{eq:beta}).

\paragraph{Aktualizacja AMAF w backpropagation.}
Dla każdego węzła $n$ na ścieżce od korzenia do liścia: niech $p_n$ oznacza
gracza na ruchu w $n$. Dla każdego ruchu $m$ zagranego w fazie in-tree
\emph{poniżej} $n$ \emph{lub} w fazie symulacji, jeśli ruch wykonał gracz
$p_n$:
\begin{itemize}[nosep]
    \item $\tilde{n}(n, m) \mathrel{{+}{=}} 1$,
    \item $\tilde{w}(n, m) \mathrel{{+}{=}} 1$, jeśli $p_n$ wygrał symulację.
\end{itemize}

Włączenie ruchów in-tree (a nie tylko symulacyjnych) jest kluczowe ---
pierwsza wersja implementacji pomijała te ruchy, prowadząc do drastycznej
porażki RAVE w eksperymentach (patrz~sekcja~\ref{sec:llm}).

\paragraph{Wybór hiperparametru $b^2$.}
Klasyczna wartość Silvera ($b^2 = 10^{-4}$) okazała się nieoptymalna
dla Breakthrough --- daje $\beta$ tak bliskie 1, że Q-RAVE dominuje selekcję.
Po zamiataniu (sweep) $b^2 \in \{10^{-6}, \ldots, 10^{-1}\}$, wybraliśmy
$b^2 = 10^{-2}$ jako najlepiej działającą wartość.

Eksponowany jako \texttt{RaveAgent(iterations, c, rave\_k, seed)},
gdzie \texttt{rave\_k} dla zachowania kompatybilności API jest interpretowane
jako $b^2$.

\subsection{UCT + Progressive Bias (Chaslot et al. 2008)}

\paragraph{Algorytm.}
Selekcja potomka rozszerza UCB1 o człon biasu malejący z liczbą odwiedzin
\cite{chaslot2008progressive}:
\begin{equation}
    \text{score}(c) = Q_{\text{UCT}}(c) + c_{\text{exp}} \sqrt{\frac{\ln N(s)}{n(c)}} + w_{\text{bias}} \cdot \frac{H(c)}{n(c) + 1}
\end{equation}
gdzie $H(c) \in [0, 1]$ jest wartością heurystyki ruchu wiodącego do $c$,
a $w_{\text{bias}}$ to waga biasu (= 1.0 w naszych eksperymentach).
Człon $1/(n(c)+1)$ powoduje naturalny zanik biasu wraz ze wzrostem
liczby symulacji --- przy rosnącej $n(c)$ statystyki Q-UCT przejmują
sterowanie selekcją.

\paragraph{Funkcja heurystyczna $H$.}
Dla ruchu z pola $f$ na pole $t$ przy graczu $p$:
\begin{equation}
H(c) = \min\!\left(1{,}0, \quad \text{adv}(t, p) + b_{\text{cap}} \cdot [\text{ruch jest biciem}]\right)
\end{equation}
gdzie $\text{adv}(t, p) = t_{\text{row}} / (R-1)$ dla białego i $(R-1-t_{\text{row}}) / (R-1)$
dla czarnego (zaawansowanie pionka znormalizowane do $[0, 1]$ względem rzędu
docelowego), a $b_{\text{cap}} = 0{,}2$ to bonus za bicie. $R$ to liczba
rzędów planszy.

Eksponowany jako \texttt{ProgressiveBiasAgent(iterations, c, bias\_weight, seed)}.

\subsection{Heurystyczny gracz alpha-beta}

\paragraph{Funkcja ewaluacji.}
$f: \text{State} \times \text{Player} \to \mathbb{R}$, dodatnia gdy korzystna
dla \emph{Player}:
\begin{equation}
f(s, p) = \pm \left( w_{\text{adv}} \cdot \Delta_{\text{advancement}} + w_{\text{mat}} \cdot \Delta_{\text{material}} + w_{\text{mob}} \cdot \Delta_{\text{mobility}} \right)
\end{equation}
gdzie:
\begin{itemize}[nosep]
    \item $\Delta_{\text{advancement}} = \sum_{\text{pionków białych}} \frac{r}{R-1} - \sum_{\text{pionków czarnych}} \frac{R-1-r}{R-1}$ ---
          różnica zaawansowania pionków (znormalizowana do rzędu),
    \item $\Delta_{\text{material}}$ = liczba pionków białych $-$ liczba pionków czarnych,
    \item $\Delta_{\text{mobility}} = |\text{legalne ruchy}|$ z znakiem (+ jeśli na ruchu jest biały, $-$ wpp).
\end{itemize}
Wagi: $w_{\text{adv}} = 1{,}0$, $w_{\text{mat}} = 3{,}0$, $w_{\text{mob}} = 0{,}1$
(dobrane na podstawie szacunków z literatury \cite{lorentz2011breakthrough}
i kalibracji wstępnej). Dla terminala: $\pm 1000$ (wartość dominująca pozostałe składniki).

\paragraph{Przeszukiwanie alpha-beta.}
Klasyczny negamax z odcięciem $\alpha\beta$, głębokość 5 półruchów.
Brak heurystyk porządkujących ruchy (move ordering) --- niepotrzebne
dla 6$\times$6 (depth 5 zajmuje $\sim$30 ms).

Eksponowany jako \texttt{HeuristicAgent(depth=5, weights=None)}.

\subsection{Architektura implementacji}

\paragraph{Stos technologiczny.}
Rdzeń obliczeniowy (silnik gry, wszystkie agenci, ewaluacja) zaimplementowany
w języku Rust. Orchestracja eksperymentów, GUI i analiza wyników w Pythonie 3.12.
Wiązanie Rust $\leftrightarrow$ Python przez PyO3 0.22 z budowaniem przez maturin.

\paragraph{Struktura repozytorium.}
\begin{verbatim}
breakthrough-mcts/
  src/                              # Rust
    game.rs                         # Bitboard GameState
    mcts/{uct,rave,progressive_bias}.rs
    heuristic.rs                    # alpha-beta
  python/breakthrough/
    reference_game.py               # Walidacja krzyżowa
    gui.py                          # Pygame GUI
    experiments/                    # Run_h{1,2,3,4}.py
  tests/test_rust_vs_python.py      # Walidacja silnika
  notebooks/                        # Skrypty analityczne
  report/                           # Niniejszy dokument
\end{verbatim}

\paragraph{Reprodukowalność.}
Generator: Xoshiro256++ \cite{vigna2017}. Master seed wyprowadza
deterministycznie seedy partii (SHA-256). Każda partia ma jawnie zapisany seed
w logu JSONL --- możliwa pełna replikacja pojedynczej partii.

\paragraph{Decyzja o zrównolegleniu.}
\textbf{MCTS jest sekwencyjny}, zrównoleglenie odbywa się wyłącznie na poziomie
batchowania eksperymentów (\texttt{ProcessPoolExecutor}). Decyzja motywowana
czystością metodologiczną porównań --- porównujemy algorytmy, nie zrównoleglenie.
Tree parallelization (virtual loss) ani root parallelization \emph{nie} są
zaimplementowane.
